\documentclass[a4paper,12pt]{report}
\renewcommand\thesection{\arabic{section}}
\renewcommand{\contentsname}{Cuprins}
\renewcommand{\figurename}{Figura}
\renewcommand{\tablename}{Tabel}
\usepackage{amsmath}
\usepackage{rotating}
\usepackage{listings}
\usepackage{lipsum}  
\usepackage[backend=biber,style=numeric]{biblatex}
\addbibresource{bibliography.bib}
\setcounter{tocdepth}{6}
\setcounter{secnumdepth}{6}

\usepackage{url}
% \usepackage[backend=bibtex,style=numeric,citestyle=numeric,backref=true,sorting=none]{biblatex}


\usepackage{indentfirst}
\usepackage[nobottomtitles*]{titlesec}

\newcommand*{\justifyheading}{\raggedright}
\titleformat{\section}{\normalfont\Large\justifyheading}{\thesection}{18pt}{}

\usepackage{graphicx}
\usepackage[export]{adjustbox}
\graphicspath{ {./images/} }

\usepackage[center]{caption}
\usepackage{subcaption}

\usepackage[margin=3.9cm]{geometry}


\usepackage{listings}
\usepackage{color}
\usepackage{float}

\definecolor{codegreen}{rgb}{0,0.6,0}
\definecolor{codegray}{rgb}{0.5,0.5,0.5}
\definecolor{codepurple}{rgb}{0.58,0,0.82}
\definecolor{backcolour}{rgb}{0.96,0.96,0.96}

\lstdefinestyle{mystyle}{
	backgroundcolor=\color{backcolour},   
	commentstyle=\color{codegreen},
	keywordstyle=\color{blue},
	numberstyle=\tiny\color{codegray},
	stringstyle=\color{codepurple},
	basicstyle=\fontsize{10}{10}\selectfont\ttfamily,
	breakatwhitespace=false,         
	breaklines=true,                 
	captionpos=b,                    
	keepspaces=true,                 
	numbers=left,                    
	numbersep=2pt,                  
	showspaces=false,                
	showstringspaces=false,
	showtabs=false,                  
	tabsize=2
}

\lstset{style=mystyle}

\usepackage{pythonhighlight}

\usepackage{fontspec}% -> Schimbat compilator in LuaLatex
\setmainfont{UTSans-Regular.otf}
\usepackage{setspace}
\usepackage{tocloft}
\setstretch{1} 

\begin{document}
	
\include{./cover/cover}
\begin{titlepage}
\begin{flushleft}
\textbf{Departamentul de Electronică și Calculatoare}\\
\textbf{Programul de studii: Calculatoare}
\end{flushleft}

 
	\begin{center}		
		\vspace{1.5cm}
		\Large \textbf{GÎRBACEA Marian Ionuț}
		
		\vspace{1cm} 
		
		\Huge Sistem interactiv de antrenare VR în mediul industrial cu roboți
		\\
		
		\vfill

        \Large
        \textbf{Conducători științifici:} \\
        conf. dr. ing.  DANCIU Gabriel Mihail\\
		şef lucr. dr. ing. BOGDAN Ioana Corina

		\vfill
		
		\Large
		Brașov, 2026
		
	\end{center}
\end{titlepage}

\newpage
\setcounter{page}{2}
% \tableofcontents

\newpage
\pagenumbering{arabic}
\renewcommand{\bibname}{Bibliografie}
\renewcommand\lstlistlistingname{\normalsize{CODURI SURSĂ}}
\renewcommand{\lstlistingname}{\normalsize{Secvență de Cod}}
\renewcommand{\listfigurename}{FIGURI}
\renewcommand{\cftloftitlefont}{\normalsize}
\renewcommand{\cftfigfont}{\normalsize}
\renewcommand{\cftfigpagefont}{\normalsize}
\renewcommand{\listtablename}{TABELE}
\renewcommand{\cftlottitlefont}{\normalsize}
\renewcommand{\cfttabfont}{\normalsize}
\renewcommand{\cfttabpagefont}{\normalsize}
\tableofcontents
% acestea sunt actualizate in mod automat, folosind table, figure si lstlisting
 \clearpage
 \section*{Lista de figuri, tabele și coduri sursă}
\addcontentsline{toc}{section}{Lista de figuri, tabele și coduri sursă}


 \begingroup
 \let\clearpage\relax
 \listoffigures
 \let\clearpage\relax
 \listoftables
 \let\clearpage\relax
 \lstlistoflistings
 %\setcounter{page}{5}
 \endgroup


 \clearpage

% % trebuie completata manual
\section*{Lista de acronime}
\addcontentsline{toc}{section}{Lista de acronime}

% (se vor scrie toate prescurtările utilizate in ordine alfabetică)

\noindent 
2D - 2 Dimensiuni;\\
3D - 3 Dimensiuni; \\
AR - Augumented Reality;\\
CAD - Computer Aided Design;\\
CAE - Computer Aided Engineering;\\
Kg - Kilograms
CAM - Computer Aided Manufacturing;\\
PPD - Pixel Per Degree;\\
PPI - Pixel Per Inch;\\
RGB - Red Green Blue;\\
%ROS - Robot Operating System;\\
HMD - Head-Mounted Display; \\
HZ - Hertz
IoT - Internet of Things;\\
MQTT - Message Queuing Telemetry Transport;\\
OASIS - Organization of Advancement of Structured Information Standard;\\
SCARA - Selective Compliance Articulated Robot Arm;\\
UI - User Interface;\\
URDF - Unified Robotics Description Format;\\
USB - Universal Serial Bus:\\
XML - Extensible Markup Language;\\
XR - Extended Reality;\\
VR - Virtual Reality;\\
WIFI - Wireless Fidelity;\\

% \clearpage
\clearpage

\section{Capitole principale }
1. Introducere
  \begin{itemize}
 
       \item  Tema proiectului.
       \item  Obiectivele și scopul lucrării. 
 \end{itemize}


2. Fundamente teoretice
  \begin{itemize}
 
       \item  Ce este tehnologia VR
       \item  Cum se poate interacționa în mediul VR?
       \item  Despre roboții industriali. 
        \item De ce este necesară simularea instruirii pe roboții industriali?
        \item Ce este URDF?
        \item Ce este MQTT?
 \end{itemize}

3. Proiectarea în VR
  \begin{itemize}
     \item Echipamente folosite(descriere)
       \item  descriere scenarii antrenament
       \item  arhitectura aplicației
       \item mecanisme de interacțiune utilizator-robot
        \item modelarea robotului/mediului pentru simulare în VR
 \end{itemize}

4. Implementare
  \begin{itemize}
 
       \item  creare mediu în VR
       \item  implementarea comportamentului robotului
       \item sistem de feedback pentru utilizator
        \item integrarea elementelor 3D
 \end{itemize}
 
 5. Testarea implementării

 6. Concluzii

\subsection{Obiective principale:}
  \begin{itemize}
       \item  modelarea în 3D a robotului.
       \item  implementare unui scenariu în VR în care utilizatorul poate interacționa cu robotul.
       \item implementare unui mod de antrenament pe diferite scenarii.
        \item evaluarea utilității sistemului ca instrument educațional.
 \end{itemize}


\section{Echipamente disponibile:}
  \subsection{Hardware:}
  \begin{itemize}
       \item  robot (MaxArm:)
       \item  cască VR (Oculus Quest 3s:)
    
 \end{itemize}

  \subsection{Software:}
  \begin{itemize}
       \item  Unity cu plugin-uri pentru dezvoltare VR.
       \item  Blender/Fusion 360 pentru modelare în 3D.
     
 \end{itemize}



\renewcommand\thesection{\arabic{section}}
\renewcommand{\contentsname}{Cuprins}
\renewcommand{\figurename}{Figura}
\renewcommand{\tablename}{Tabel}

\setcounter{tocdepth}{6}
\setcounter{secnumdepth}{6}


% \usepackage[backend=bibtex,style=numeric,citestyle=numeric,backref=true,sorting=none]{biblatex}




\titleformat{\section}{\normalfont\Large\filleft}{\thesection}{18pt}{}



\definecolor{codegreen}{rgb}{0,0.6,0}
\definecolor{codegray}{rgb}{0.5,0.5,0.5}
\definecolor{codepurple}{rgb}{0.58,0,0.82}
\definecolor{backcolour}{rgb}{0.96,0.96,0.96}

\lstdefinestyle{mystyle}{
	backgroundcolor=\color{backcolour},   
	commentstyle=\color{codegreen},
	keywordstyle=\color{blue},
	numberstyle=\tiny\color{codegray},
	stringstyle=\color{codepurple},
	basicstyle=\fontsize{10}{10}\selectfont\ttfamily,
	breakatwhitespace=false,         
	breaklines=true,                 
	captionpos=b,                    
	keepspaces=true,                 
	numbers=left,                    
	numbersep=2pt,                  
	showspaces=false,                
	showstringspaces=false,
	showtabs=false,                  
	tabsize=2
}

\lstset{style=mystyle}




\setmainfont{UTSans-Regular.otf}


	


\newpage



\newpage
\pagenumbering{arabic}
\renewcommand{\bibname}{Bibliografie}
\renewcommand\lstlistlistingname{\normalsize{CODURI SURSĂ}}
\renewcommand{\lstlistingname}{\normalsize{Secvență de Cod}}
\renewcommand{\listfigurename}{FIGURI}
\renewcommand{\cftloftitlefont}{\normalsize}
\renewcommand{\cftfigfont}{\normalsize}
\renewcommand{\cftfigpagefont}{\normalsize}
\renewcommand{\listtablename}{TABELE}
\renewcommand{\cftlottitlefont}{\normalsize}
\renewcommand{\cfttabfont}{\normalsize}
\renewcommand{\cfttabpagefont}{\normalsize}
 \clearpage

% % trebuie completata manual

\clearpage




 \setcounter{page}{4}
\section{Context și motivație}

\subsection{Context}

În societatea contemporană majoritatea domeniilor cum sunt cele de producție, farmaceutică până la apărare sau alimentar prezența echipamentelor de tip roboți sau brațe robotice a devenit tot mai răspândită, unde se  integrează expertiza umane cu cea a asistenței din partea roboților. Astfel a apărut o nevoie pentru testarea și găsirea mai rapidă a erorilor și defectelor unor astfel de echipamente din partea producătorilor și antrenarea directă pe aceștia poate duce la defectarea lor sau chiar la accidentări, ducând la costuri suplimentare atât din punct de vedere al companiei care a achiziționat robotul cât și din partea producătorului robotului.

O soluție pentru testarea și găsirea eventualelor erori și defecte o reprezintă tehnologia "digital twin". Pentru partea de antrenare a operatorului uman se poate folosi realitatea virtuală pentru a implementa un mediu de antrenament ce implementează un "digital twin", mediul cuprinzând majoritatea scenariilor în care operatorul uman și brațul robotic pot să fie supuse în activitățile la care sunt implicate.
%cuprizând toate funcționalitățile brațului real și nefiind necesară interacțiunea directă cu brațul. 


O interacțiune directă dintre robotul real și "digital twin" poate fi utilă în momentul în care este necesară controlul direct al brațului real sau al robotului în medii periculoase cum sunt cele manevrarea elementelor radioactive sau a fiolelor din laboratoare.

Lucrarea își propune realizarea unui astfel de mediu de antrenament în realitatea virtuală folosind implementarea unui  "digital twin"  pentru un braț robotic cu 4 grade de liberate ce permite ridicarea și mutarea de obiecte folosind o pompă. De asemenea acest sistem urmărește ca procesul de antrenare a utilizatorilor umani, care nu au interacționat niciodată sau rar cu un astfel de sistem să fie cât mai simplă, intuitivă astfel încât timpul petrecut în antrenare să fie cât mai mic și nivelul de pregătire dobândit să fie cât mai mare.


%sunt tot mai prezente dispozitive și echipamente electronice care ne ajută și ne îmbunătățesc stilul de viață în majoritatea privințelor. Astfel au apărut 2 probleme principale:
%aceste echipamente și dispozitive sunt construite din mii și mii de piese și componenete și aceste echipamente și dispozitive a apărut o cerere din ce în ce mai mare. 

%Având în vedere aceste aspecte, cererea și complexitatea, dexteritatea și precizia umamă a ajuns să nu mai fie suficientă. Astfel pentru procesele de fabricație au început să fie folosite aproape în totalitate, roboți și brațe robotice, care fie sunt controlate în mod automat, fie sunt controlate de către un operator uman, apărând nevoia antrenării unui om despre cum să folosească aceste brațe/roboți

%În ultimul timp colaborarea dintre oameni și roboți în mediile industriale a prins tracțiune, în special în procesele de asamblare și dezasamblare de mare precizie. Creșterea complexității mediului de lucru industrial, necesită metodologii de antrenare avansate, prin integrarea expertizei umane cu cea a asistenței din partea roboților. Dar această instruire necesită resurse ridicate și implică riscuri de accidentare.


%Tehnologia VR 

%Această lucrare își propune realizarea unui sistem interactiv de antrenare folosind realitatea virtuală pentru mediul industrial care lucrează cu roboți care să faciliteze antrenarea utilizatorilor umani neexperimentați la un mediul complet de antrenare în folosirea roboților fără implicarea de riscuri.

\subsection{Motivație}
Motivația principală din spatele realizării constă în creșterea în popularitate a tehnologiei de "digital twin" în domenii precum:medical, de producție  domeniul energetic \cite{DigTwin}.
Pe partea de producție aceștia pot ajuta operatorul să înțeleagă condițiile asset-ului și să realizeze planuri predictive pentru mentenanță. Producătorii folosind "digital twins" pentru planificarea zonei de producție, pot să simuleze impactul asupra liniei de producție, ajustând și optimizând pentru cel mai bun flux de lucru. În domeniul auto, producătorii deja folosesc "digital twins" pentru a prototipa și simula modele noi de vehicule. Aceștia pot să ofere posibilitatea pentru vehicule pe bază de software și mașini autonome \cite{DT}.
Pentru domeniul energetic, "digital twins" pot fi folosiți în rețele electrice și tehnologia microgrid, pentru a ajuta operatorii în optimizarea și distribuția de energie. Flexibilitatea acestei tehnologii le oferă operatorilor noi unelte pentru a răspunde într-un mod proactiv la deranjamente, să simuleze situații de urgență și în plănuirea integrării surselor de energie regenerabilă eficientă din punct de vedere al costului.
În medicină, le oferă furnizorilor de servicii medicale posibilitatea îmbunătățirii uneltelor de monitorizarea a pacienților. Prin combinarea datelor de la dispozitive multiple destinate pacienților într-un "digital twin", furnizorii pot să obțină un o vedere mai mare de asamblu asupra sănătății pacientului. De asemenea, pot ajuta administratorii să monitorizeze rata de ocupare și îngrijorările legate de siguranță sau chiar să monitorizeze calitatea aerului pentru a ajuta în prevenirea răspândirii patogenilor \cite{DT}.

Un alt domeniu în care această tehnologie poate fi folosită este cea a telecomunicațiilor. Aceasta le poate oferi furnizorilor din telecom să obține o imagine de ansamblu mai bună asupra  infrastructurii rețelei, a stării asset-ului și să ofere posibilitatea de a se realiza mentenanță predictivă pentru minimizarea momentelor de picare a rețelei. Pe lângă furnizorii îmi mai pot folosi în simularea fluctuării cererilor din rețea pentru teste de stres sau planificarea alternativelor fără necesitate de schimbări costisitoare în lumea reală \cite{DT}.



Un avantaj principal al acestui proiect constă în folosirea tehnologie de realitate virtuală, care reprezintă o soluție modernă care permite simularea și controlul acestor roboți, oferindu-le utilizatorilor posibilitatea de a învăța cum să-i controleze, limitele lor și scenariile de utilizare fără a fi necesar interacțiunea fizică cu cei reali. Un sistem în realitatea virtuală poate reduce costurile, crește eficiența în antrenare și de a îmbunătăți nivelul de pregătire fără a afecta productivitatea reală.

Un alt aspect important al proiectului constă în implementarea a 2 metode de control al brațului din VR. Această metodă de implementare a metodei de control prezintă un avantaj semnificativ deoarece oferă utilizatorului posibilitatea de control al robotului într-un mod simplu și ușor de înțeles, cât și producătorului să testeze rapid implementări și să rezolve eventualele probleme apărute.





\section{Obiectivele lucrării}

Principalele obiective stabilite pentru realizarea lucrării:
\begin{itemize}
	\item Dezvoltatea unui sistem de antrenament care să ofere un minim de pregătire asupra controlului robotului
	\item Realizarea modelului 3D al brațului din realitate
	\item Importarea modelului 3D al brațului în mediul VR din Unity
	\item Aplicarea proprietăților și constrângerilor robotului care sunt necesare pentru mișcarea acestuia
	\item Realizarea primului mod de mișcare al robotului baza pe asociere directă între butoanele de la manete și motoare
	\item Realizarea celui de al 2-lea mod de control al robotului bazat pe cinematică inversă
	\item Aplicarea valorilor citite din Unity asupra robotului din lumea reală pentru a putea fi controlat
	\item Primirea feedback-ului de la robotul real legat de valorile pe care le-a recepționat
\end{itemize}


\clearpage
%Scopul lucrării este de a crea un mediu de simulare interactivă , simplu și sigur pentru antrenarea în folosirea roboților care este destinată utilizatorilor neexperimentați.

%De asemenea acest sistem urmărește ca procesul de antrenare a utilizatorilor umani, care nu au interacționat niciodată sau rar cu un astfel de sistem să fie cât mai simplă, intuitivă astfel încât timpul petrecut în antrenare să fie cât mai mic și nivelul de pregătire dobândit să fie cât mai mare.

%Nu în ultimul rând prin folosirea acestui mediu de antrenare procesul de antrenare nu implică riscurile aferente, dacă antrenarea implica roboții reali.



%Această lucrare prezintă următoarea structură:capitolul 2 prezintă tehnologiile și fundamentele teoretice care stau la baza realizării acestui proiect, capitolul 3 prezintă echipamentele folosite, implementarea, arhitectura aplicației în  realitatea virtuală pentru și modul în care robotul a fost transpus în acest mediu, iar capitolul 4 cuprinde testele realizate și rezultatele în urma testelor, modul în care au fost îndeplinite obiectivele, utilitatea acestui proiect și direcțiile viitoare.
     



% \nocite{*}




\renewcommand\thesection{\arabic{section}}
\renewcommand{\contentsname}{Cuprins}
\renewcommand{\figurename}{Figura}
\renewcommand{\tablename}{Tabel}

\setcounter{tocdepth}{6}
\setcounter{secnumdepth}{6}


% \usepackage[backend=bibtex,style=numeric,citestyle=numeric,backref=true,sorting=none]{biblatex}




\titleformat{\section}{\normalfont\Large\justifyheading}{\thesection}{18pt}{}



\definecolor{codegreen}{rgb}{0,0.6,0}
\definecolor{codegray}{rgb}{0.5,0.5,0.5}
\definecolor{codepurple}{rgb}{0.58,0,0.82}
\definecolor{backcolour}{rgb}{0.96,0.96,0.96}

\lstdefinestyle{mystyle}{
	backgroundcolor=\color{backcolour},   
	commentstyle=\color{codegreen},
	keywordstyle=\color{blue},
	numberstyle=\tiny\color{codegray},
	stringstyle=\color{codepurple},
	basicstyle=\fontsize{10}{10}\selectfont\ttfamily,
	breakatwhitespace=false,         
	breaklines=true,                 
	captionpos=b,                    
	keepspaces=true,                 
	numbers=left,                    
	numbersep=2pt,                  
	showspaces=false,                
	showstringspaces=false,
	showtabs=false,                  
	tabsize=2
}

\lstset{style=mystyle}




\setmainfont{UTSans-Regular.otf}





\newpage



\newpage
\pagenumbering{arabic}
\renewcommand{\bibname}{Bibliografie}
\renewcommand\lstlistlistingname{\normalsize{CODURI SURSĂ}}
\renewcommand{\lstlistingname}{\normalsize{Secvență de Cod}}
\renewcommand{\listfigurename}{FIGURI}
\renewcommand{\cftloftitlefont}{\normalsize}
\renewcommand{\cftfigfont}{\normalsize}
\renewcommand{\cftfigpagefont}{\normalsize}
\renewcommand{\listtablename}{TABELE}
\renewcommand{\cftlottitlefont}{\normalsize}
\renewcommand{\cfttabfont}{\normalsize}
\renewcommand{\cfttabpagefont}{\normalsize}
\clearpage

% % trebuie completata manual

\clearpage




\setcounter{page}{7}
\section{Fundamente teoretice}
\subsection{Digital Twin}
Un "digital twin" este o reprezentare virtuală a unui obiect fizic sau sistem ce folosește date din timp real să reflecte cu acuratețe comportamentului geamănului din lumea reală, performanțe și condiții.
"Digital twins" permit monitorizate continuă, simulare și analiză a unui obiect, sistem sau produs pe durata timpului său de viață, de la proiectare și producție, până la mentenanță și dezafectare.
Aceștia pot  să reprezinte orice obiect într-un mod virtual, de la clădiri și poduri, până la mașini, avioane, artefacte antice și chiar întreaga planetă\cite{DigTwinIBM}. 

De asemenea, prin "digital twins"  pot să se modeleze sisteme complexe cum sunt: traficul, evenimente meteorologice, planuri medicale de tratament, procese și operațiuni din fabrici sau în medii experimentale, și chiar pot să fie bazați fie pe oameni reali, fie imaginari cu toate trăsăturile necesare cum sunt: tonul vocii, înfățișare și personalitate \cite{DigTwinIBM}. 


\subsection{Cinematica inversă}
Cinematica reprezintă studiul mișcării fără să se considere cauza mișcării, cum sunt forțele și cuplurile. Cinematica inversă reprezintă folosirea ecuațiilor cinematice să determine mișcarea unui robot pentru a ajunge la destinație. De exemplu, pentru a putea apuca un obiect, un braț robotic folosit în procese de fabricație necesită mișcări precise de la poziția inițială către poziția necesară. Capătul de prindere al robotului reprezintă efectorul final. Configurația robotului reprezintă o listă a pozițiilor de "Joint-uri" care pot fi acționate atât timp cât nu depășesc constrângerile robotului.

 Pornind de la poziția la care efectorul final trebuie să ajungă, cinematica inversă poate să determine o configurație a  "Joint-urilor" astfel încât efectorul final să poată să se miște către țintă \cite{IK}.

\subsection{Realitatea virtuală}
Realitatea virtuală reprezintă un termen care descrie un mediu 3D generat de un calculator care poate fi explorat și se poate interacționa de către un utilizator uman. Acel utilizator devine o parte a mediului virtual sau devine imersive cu mediul și cât timp se află în mediu, poate să manipuleze obiectele sau să performeze o serie de acțiuni.

Realitatea virtuală este implementată folosind sisteme care sunt destinate pentru realizarea mediului, cum sunt căștile, benzile de mers multi-direcționale și mănuși speciale. Aceste sisteme sunt folosite să stimuleze simțurile simultan pentru a crea o iluzie a realității

% un mediu virtual în 3 dimensiuni, generat de calculator în care utilizatorul poate interacționa cu mediul. Accesul la acest mediu este realizat prin intermediul unui calculator care este capabil să afișeze informație 3D prin intermediul unui ecran, care poate fi fie ecrane izolate sau un ecran care poate care pot fi prins de capul utilizatorului împreună cu o serie de senzori folosiți\cite{vr}. 

Această tehnologie poate fi împărțită în 5 tehnologii principale: 
\begin{itemize}
	\item Realitatea virtuală non-imersivă se referă la utilizarea unor serii de ecrane care înconjoară utilizatorul pentru ai fi prezentată informația virtuală\cite{vr}.
	
	\item Realitatea virtuală semi-imersivă în care componentele digitale se suprapun peste obiectele reale. Astfel aceste elemente virtuale pot fi folosite într-un mod similar cu cele reale \cite{vr3}.
	
	\item Realitatea virtuală imersivă, care a fost folosită în acest proiect se folosește de folosirea unui HMD pentru a urmării mișcările utilizatorului și să prezinte informația VR în funcție de mișcările utilizatorului\cite{vr}.
	
	\item Realitatea augmentată care reprezintă un alt mod de a realiza realitate virtuală prin modul în care dezvoltatorii suprapun cele ambele realități \cite{vr3}.
	
	\item Realitate mixtă combină lumea fizică cu cea virtuală. Reprezintă un caz special de realitate augmentată. Această tehnologie permite vizualizarea de oameni și obiecte într-un context real \cite{vr3}.
\end{itemize}






\subsection{MQTT}

MQTT reprezintă un standard  OASIS fiind un protocol de mesagerie folosit în IoT. Este proiectat ca să implementeze protocolul de tip publicare/abonare într-un mod în care să foloseacă cât mai puține resurse. Acest protocol este ideal pentru conectarea între dispozitive, ce sunt limitate din punct de vedere hardware. Comunicarea prin acest protocol este bidirecțională. \cite{MQTT}.
%ROS reprezintă un set de biblioteci open source software și unelte care sunt folosite în aplicațiile care implică roboți în special controlul componentelor roboților de la un calculator \cite{ros}.

%Un sistem ROS este format dintr-o serie de noduri independente, care comunică între ele folosind un model de mesagerie de tip publicare/abonare \cite{ros}. 

Mesageria de tip publicare/abonare reprezintă un model de comunicare asincron care permite dezvoltatorilor să construiască aplicații funcționale și complexe în cloud. Pentru aceste sisteme cloud modelul de mesagerie oferă notificări instantanee când se întâmplă un eveniment în sistemele distribuite \cite{ps}.

Acest model funcționează în felul următor:
\begin{itemize}
	\item Mesajele reprezintă un set de date de comunicație care sunt trimise de la emițător la receptor. Mesajele pot fi de orice tip de dată \cite{ps}
	\item Topicul care acționează ca un canal intermediar între emițător și receptor. Gestionează o listă de receptori care sunt interesate în mesajele care au același topic \cite{ps}
	\item Abonații reprezintă destinatarul mesajului. Aceștia trebuie să se înregistreze sau să se aboneze la un topicul care îi interesează. Pot să performeze diferite funcții sau să facă ceva diferit cu mesajul în paralel \cite{ps}
	\item Publicatorii reprezintă componenta care trimite mesajele. Creează mesaje despre un topic și le trimite pe toate la toți abonații ai acelui topic. Interacțiunea dintre publicator și abonați reprezintă o relație de tipul unul la mulți. Publicatorul nu trebuie să știe cine folosește informația, iar abonații nu trebuie să cunoască proveniența mesajelor, transmisia fiind una de difuzare \cite{ps}
\end{itemize}

Modelul de mesagerie de tip publicare/abonare reprezintă următoarele avantaje:
\begin{itemize}
	\item Eliminarea polling-ului. Topicurile de mesaje care permit transmiterea mesajelor instantaneu eliminând verificările periodice pe care consumatorii mesajului trebuie să le facă când apare informație nouă sau când este actualizată \cite{ps}.
	\item Direcționare dinamică. Modelul plubicare/abonat face ca descoperirea serviciilor mai ușor, mai natural, și mai greu susceptibilă la erori. În loc ca aplicația să întrețină un grup de abonați căruia să-i trimită mesajele, un publicator va posta mesajele către un topic specific, astfel încât abonații să se aboneze la mesajele cu topic-urile de interes \cite{ps}.
	\item Simplificarea comunicației. Acest model simplifică procesul de comunicație prin eliminarea conexiunilor punct-la-punct care leagă o singură conexiune către un topic. Topicul va gestiona abonările, astfel încât să decidă ce mesaje sunt destinate cărui endpoint \cite{ps}.
\end{itemize}

Aceste mesaje pot fi consumate de oricare număr de noduri, incluzând, filtre, sisteme de nivel înalt cum ar fi cele de ghidare \cite{ros}. 

%ROS nu trebuie să fie pe același sistem sau pe aceeași arhitectură. Pot exista simultan componente diferite care să publice mesaje, să se aboneze la mesaje și altele care să controleze componentele. Acestea fac ca ROS să fie flexibil șă să se adapteze la necesitățule aplicației \cite{ros}.

\subsection{URDF}
URDF este un limbaj markup bazat pe XML, dezvoltat ca parte a ecosistemului ROS. Este folosit în a reprezenta modelul unui robot în raport cu: structura fizică, cinematica, proprietățile legate de inerție, aspectul fizual și geometria coliziunilor \cite{urdf}.

Un fișier definește un robot ca o colecție de legături (corpuri rigide), conectate prin articulații (care definește cum legăturile se pot mișca între ele). De asemenea poate specifica senzori, transmisii și materiale, permițând simularea și vizualizarea robotului în mediul 3D \cite{urdf}.

URDF servește ca o componentă esențială în controlul prin:
\begin{itemize}
	\item Simulare: Modelele URDF sunt folosite pentru a încărca roboții în simulatoare cum este Gazebo, sau medii de dezvoltare 3D cum este Unity, unde dezvoltatorii testează algoritmi în medii de dezvoltare virtuale controlate înainte să fie încărcate pe hardware. URDF se asigură că constrângerile fizice și reprezentările vizuale se aseamănă cu cele din lumea reală \cite{urdf}.
	\item Cinematica și planificarea mișcărilor. Este necesară pentru a ajuta la înțelegerea limitelor articulațiilor și ierarhiei dintre legături și coordonatele componentelor. Acestea permit planificarea mișcărilor cu acuratețe și calculului cinematicii inverse și directe \cite{urdf}.
	\item Vizualizare și Debbuging. URDF este folosit pentru a randa vizualizări 3D ale robotului în timp real. Se pot suprapune datele de la senzori pe model, ceea ce face mai ușor înțelegerea comportamentului sistemului și debugging-ul \cite{urdf}.
	\item Integrarea cu senzori. URDF permite poziționarea cu precizie a cadrelor și a pozițiilor senzorilor pe robot. Modelarea acestora cu acuratețe este esențială pentru combinarea senzorilor și a navigației \cite{urdf}.
\end{itemize}


\subsection{Roboții industriali}
Roboții industriali sunt folosiți în procese de asamblare, curățare și celelalte activități din fabrici unde oamenii nu pot opera. Automatizarea fiecărui proces îmbunătățește eficiența și contribuie la reducerea lipsei forței de muncă și al costurilor \cite{ri}.

Roboții industriali se împart în 6 tipuri:
\begin{itemize}
	\item Roboți liniari. Aceștia sunt asemănători în  aspect cu o macara de tip pod cu 3 axe de mișcare, are o capacitate de mutare de obiecte și acuratețe ridicată. Axele acestui robot sunt perpendiculare între ele și robotul se mișcă cu axele în linie dreaptă \cite{ri1}.
	\item Roboți cilindrici. Aceștia sunt roboți pe 3 axe cu zone de lucru sub formă de cilindri, 2 axe ale robotului se mișcă linear, iar cel de al 3 lea se mișcă circular \cite{ri1}.
	\item Roboți sferici. Aceștia au sistem de coordonate polare. Au 2 axe de rotație și o axă lineară.
	\item Roboți SCARA. Sunt roboți de mare precizie și viteză ce au 4 axe de mișcare. Încheieturile roboților se mișcă pe un plan orizontal împreună cu alte 3 axe, și o altă componentă se mișcă sus și jos împreună cu cea de a 4 a axă \cite{ri1}.
	\item Roboți articulați. Sunt roboți pe 6 axe compuse din încheieturi conectate secvențial, care copiază aspectul unei mâini umane. Cele 6 axe oferă robotului flexibilitate și le permite axelor să ajungă în zone în care ceilalți nu pot \cite{ri1}.
	\item Robot delta. Sunt roboți ce au 3 mânere atașate de o bază și care țin componenta de lucru a robotului paralelă cu baza. Aceștia pot avea 3, 4 și 6 axe de mișcare. Varianta cu 3 axe poate muta obiectele pe planuri paralele, și variantele cu un număr mai mare de axe sunt capabil să rotească obiecte și să manipuleze obiecte din părți diferite \cite{ri1}.
\end{itemize}


% URDF reprezintă un fișier de tip XML care include descrierea fizică a robotului. Este în principiu un model 3D ce conține informații despre încheieturile robotului, motoare, greutate. 

% Fișierele sunt folosite în ROS, iar informațiile acestora informează operatorul uman despre aspectul, dimensiunea și capabilitățile robotului și capabilitățile acestuia înainte să înceapă să folosească robotul.

% Pe lângă cele mai de sus fișierele URDF mai sunt folosite în testarea comportamentală a roboților, până să fie puși în funcționare, ori pentru crearea de dubluri digitale folosirea în virtualizarea stării în timp real al robotului.


% \nocite{*}




\renewcommand\thesection{\arabic{section}}
\renewcommand{\contentsname}{Cuprins}
\renewcommand{\figurename}{Figura}
\renewcommand{\tablename}{Tabel}

\setcounter{tocdepth}{6}
\setcounter{secnumdepth}{6}


% \usepackage[backend=bibtex,style=numeric,citestyle=numeric,backref=true,sorting=none]{biblatex}




\titleformat{\section}{\normalfont\Large\justifyheading}{\thesection}{18pt}{}



\definecolor{codegreen}{rgb}{0,0.6,0}
\definecolor{codegray}{rgb}{0.5,0.5,0.5}
\definecolor{codepurple}{rgb}{0.58,0,0.82}
\definecolor{backcolour}{rgb}{0.96,0.96,0.96}

\lstdefinestyle{mystyle}{
	backgroundcolor=\color{backcolour},   
	commentstyle=\color{codegreen},
	keywordstyle=\color{blue},
	numberstyle=\tiny\color{codegray},
	stringstyle=\color{codepurple},
	basicstyle=\fontsize{10}{10}\selectfont\ttfamily,
	breakatwhitespace=false,         
	breaklines=true,                 
	captionpos=b,                    
	keepspaces=true,                 
	numbers=left,                    
	numbersep=2pt,                  
	showspaces=false,                
	showstringspaces=false,
	showtabs=false,                  
	tabsize=2
}

\lstset{style=mystyle}




\setmainfont{UTSans-Regular.otf}





\newpage



\newpage
\pagenumbering{arabic}
\renewcommand{\bibname}{Bibliografie}
\renewcommand\lstlistlistingname{\normalsize{CODURI SURSĂ}}
\renewcommand{\lstlistingname}{\normalsize{Secvență de Cod}}
\renewcommand{\listfigurename}{FIGURI}
\renewcommand{\cftloftitlefont}{\normalsize}
\renewcommand{\cftfigfont}{\normalsize}
\renewcommand{\cftfigpagefont}{\normalsize}
\renewcommand{\listtablename}{TABELE}
\renewcommand{\cftlottitlefont}{\normalsize}
\renewcommand{\cfttabfont}{\normalsize}
\renewcommand{\cfttabpagefont}{\normalsize}
\clearpage

% % trebuie completata manual





\setcounter{page}{9}

\section{Tehnologii folosite}
\subsection{Unity}
Unity reprezintă o platformă interactivă de creare de conținut 3D,  2D pentru VR și AR. \cite{unity}.

Acesta este un motor de creare de conținut în timp real pentru crearea de experiențe care sunt interactive într-un mod complet care este destinat atât utilizatorului final cât și pentru creatorii de conținut \cite{unity}.

Termenul de timp real descrie rapiditatea cu care o imagine este randată sau afișată pe un ecran chiar și când aceasta este schimbată de către utilizator. Scopul programelor software de timp real este de a randa imaginile atât de rapid încât interacțiunea cu mediul virtual să se facă fără observarea acută a întârzierilor\cite{unity}.

Unity are aplicații în următoarele domenii:media, divertisment, arhitectură
inginerie, construcții, automotiv, transport și producție\cite{unity}.

Unity în acest proiect a fost folosit pentru simplitatea cu care acest motor de creare a conținutului 3D integrează tehnologii de modelare de conținut 3D, cât și tehnologii de vizualizare și interacțiune cu roboții\cite{unity}.

\subsubsection{Unity XR Interaction Toolkit}
XR Interaction Toolkit este un sistem de nivel înalt și bazat pe componente de interacțiune pentru crearea de aplicații VR și AR. Oferă un framework care poate fi folosit în dezvoltarea de interacțiuni de 3D și UI, acestea fiind disponibile din Unity Input Events. Centrul acestui sistem este bazat pe un interactor și componente interacționabile, și un gestionator de interacțiuni care leagă împreună aceste tipuri de componente. De asemenea conține componente care pot fi folosite pentru locomoție și desenare \cite{XRI}.

XR Interaction ToolKit conține următorul set de componente care suportă următoarele acțiuni de interacțiune:
\begin{itemize}
	\item Compatibilitate cu o gamă înalte de platforme folosite pentru control: Meta Quest, OpenXR, Windows Mixed Reality și altele \cite{XRI}
	\item Manipulare de bază a obiectelelor din scenă \cite{XRI}
	\item Feedback haptic prin manetele de XR \cite{XRI}
	\item Feedback vizual pentru indicarea interacțiunilor active și posibile \cite{XRI}
	\item Interacțiune utilă cu XR Origin, ce reprezintă o componentă de cameră VR folosită în manevrarea staționară și în funcție de mărimea zonei a experiențelor VR \cite{XRI}
\end{itemize}

\subsubsection{Articulation Body}
Articulation body reprezintă o componentă ce permite construirea de articulații fizice cum sunt brațele robotice, lanțuri cinematice cu GameObjects care sunt organizate într-un mod ierarhic de tipul părinte-copil sau arbore. Acestea ajută în obținerea de fizică comportamentală realistică în contextul simulărilor și aplicațiilor industriale \cite{ABC}.



Un articulation body permite definirea pentru un singur obiect proprietăți, care pot fi definite de alte componente asemănătoare cum sunt RigidBody sau regular Joint într-o configurație clasică. Aceste proprietăți depind de tipul de poziția obiectului în ierarhie \cite{ABC}. 

Acestea sunt:
\begin{itemize}
	\item pentru obiectul rădăcină sau baza:masa, mutabilitatea, și dacă să folosească gravitația \cite{ABC}
	\item pentru obiectele copil sau brațe: masa, dacă să folosească gravitația, tipul pe rotație al ariticulației, legarea de ancora părintelui, poziția și rotația ancorei, fricțiunea articulației, amortizarea liniară și unghiulară, gradele de libertate ale mișcării, limitele unghiulare superioare și inferioare, rigiditatea, amortizarea, limita forței, poziția unghiului și viteza de deplasare \cite{ABC}.
\end{itemize}
Articulațiie din articulation body sunt rezolvate folosind metoda Featherstone care lucrează folosind coordonare reduse prin: fiecare corp, copil sau componentă are coordonare relative față de  părintele acestuia dar numai în jurul gradelor de libertate pe care le are setate. Acestea garantează că nu vor exista întinderi care nu sunt necesare \cite{ABC1}.

\subsubsection{Unity URDF Importer}
Unity URDF Importer reprezintă un pachet de Unity ce permite importarea unui robot definit prin formatul URDF într-o scenă din Unity. Importatorul pasează un fișier URDF și îl importă în Unity folosind articulation body \cite{UU}.

\subsection{Fusion 360}
Fusion 360 reprezintă o platformă care combină CAD, CAM, CAE, PCB, gestionarea datelor într-o singură și integrată platformă software în cloud \cite{Fusion}.

Fusin este capabil să:
\begin{itemize}
	\item Să explorere iterații ale proiectului cu ușurință folosind unelte de modelare 3D \cite{Fusion}
	\item Să producă părți de mașini CNC la o calitate înaltă cu uneltele de CAM/CAM integrate \cite{Fusion}
	\item Să se acceseze proiecte electronice întregi \cite{Fusion}
	\item Oferă unelte de simulare 3D pentru proiect \cite{Fusion}
	\item Se poate explora rezultate care sunt pregătite pentru producție cu proiect generativ  \cite{Fusion}
\end{itemize}

\section{Echipamente folosite}

\subsection{MaxArm}

MaxArm este un brat robotic dezvoltat de compania Hiwonder. Este un brat robotic open-source care este controlat de un microcontroler ESP32. Acesta este un robot care poate fi folosit în scop educativ sau în proiecte avansate. Robotul este echipat cu servo motoare de înaltă calitate și o pompă de subțiune. 

Are 4 grade de libertate:stânga-dreapta, sus-jos, față-spate, și rotația de obiecte. Prin aceste grade și tehnologia de cinematică inversă, acesta poate să execute o gamă largă de sarcini cum ar fi: apucare, sortare, transportare și stivuire de obiecte. Acesta suportă comunicare prin Bluetooth și WIFI și programare folosind Python prin Micro-Python și Arduino prin C/C++ \cite{MaxArm}. 

Specificații \cite{MaxArm}:
\begin{itemize}
	\item Dimensiuni: 158*160*260mm (lungime*lățime*înălțime)
	\item Greutate: 1.3Kg
	\item Metal și placă din fibră de sticlă
	\item Grade de libertate: 4
	\item Alimentare: 12V 5A pe curent continuu
	\item Controler: ESP32
	\item Metode de comunicație: WIFI și Bluetooth
	\item Motoare: HTS-35H și LFD-01M 
	\end{itemize}
	

\begin{figure}[h]
	\centerline{\includegraphics[width=\columnwidth,keepaspectratio]{MaxArm.png}}
	\caption[(Braul Robotic MaxArm)]{Brațul Robotic MaxArm}
	\label{fig:top}
\end{figure}


\clearpage
\subsection{Oculus Quest 3S}

Oculus Quest 3S sunt o pereche de ochelari de realitate mixtă de buget dezvoltată de compania Meta \cite{Quest3S}. 

Specificații \cite{Quest3S}:
\begin{itemize}
	\item Procesor: Snapdragon XR2 Gen 2
	\item Camere: 2 camere RGB cu PPD de 18 
	\item Stocare: 128GB/256GB
	\item DRAM: 8GB
	\item Rezoluție ecran: 1832x1920 pixeli/ochi cu 20 PPD și 773 PPI
	\item Rată de reîmprospătare: 72Hz, 90Hz, 120Hz
	\item Unghi de vizualizare: 96 de grade orizontal, 90 de grade vertical
	\item Optică: lentile fresnel
	\item Conectivitate: USB-c și WIFI
	\item Durată baterie: aproximativ 2.6 ore
	\item Timp de încărcare: aproximativ 1.9 ore pe 18Watts

\end{itemize}

\begin{figure}[h]
	\centerline{
	\includegraphics[width=0.6\textwidth]{Quest3S.png}}
	\caption[(Quest3S)]{Ochelari de realitate virutală Oculus Quest3S}
	\label{fig:top}
\end{figure}
% \nocite{*}





\renewcommand\thesection{\arabic{section}}
\renewcommand{\contentsname}{Cuprins}
\renewcommand{\figurename}{Figura}
\renewcommand{\tablename}{Tabel}

\setcounter{tocdepth}{6}
\setcounter{secnumdepth}{6}


% \usepackage[backend=bibtex,style=numeric,citestyle=numeric,backref=true,sorting=none]{biblatex}




\titleformat{\section}{\normalfont\Large\justifyheading}{\thesection}{18pt}{}



\definecolor{codegreen}{rgb}{0,0.6,0}
\definecolor{codegray}{rgb}{0.5,0.5,0.5}
\definecolor{codepurple}{rgb}{0.58,0,0.82}
\definecolor{backcolour}{rgb}{0.96,0.96,0.96}

\lstdefinestyle{mystyle}{
	backgroundcolor=\color{backcolour},   
	commentstyle=\color{codegreen},
	keywordstyle=\color{blue},
	numberstyle=\tiny\color{codegray},
	stringstyle=\color{codepurple},
	basicstyle=\fontsize{10}{10}\selectfont\ttfamily,
	breakatwhitespace=false,         
	breaklines=true,                 
	captionpos=b,                    
	keepspaces=true,                 
	numbers=left,                    
	numbersep=2pt,                  
	showspaces=false,                
	showstringspaces=false,
	showtabs=false,                  
	tabsize=2
}

\lstset{style=mystyle}




\setmainfont{UTSans-Regular.otf}





\newpage



\newpage
\pagenumbering{arabic}
\renewcommand{\bibname}{Bibliografie}
\renewcommand\lstlistlistingname{\normalsize{CODURI SURSĂ}}
\renewcommand{\lstlistingname}{\normalsize{Secvență de Cod}}
\renewcommand{\listfigurename}{FIGURI}
\renewcommand{\cftloftitlefont}{\normalsize}
\renewcommand{\cftfigfont}{\normalsize}
\renewcommand{\cftfigpagefont}{\normalsize}
\renewcommand{\listtablename}{TABELE}
\renewcommand{\cftlottitlefont}{\normalsize}
\renewcommand{\cfttabfont}{\normalsize}
\renewcommand{\cfttabpagefont}{\normalsize}
\clearpage

% % trebuie completata manual





\setcounter{page}{18}

\begin{comment}
	
	
	\section{Pașii de implementări}
	Pentru a putea crea un geamăn digital pentru brațul robotic a trebuit realizate următorii pași:
	
	\begin{itemize}
		\item Pentru primul pas a fost necesar crearea modelului 3D al brațului
		\item Al doilea pas a fost exportarea modelului 3D din Fusion în format de fișier URDF și importare în Unity
		\item Al treilea pas după impotare a fost setarea valorilor de amortizare, frecare, și limitele superioare și inferioare.
		\item Al patrulea pas a constat în aplicarea de constrângeri între componente în jurul motoarelor.
		\item Al cincilea pas după constrângeri a fost crearea sistemului de control
		\item Al saselea pas a constat în transmiterea datelor provenite de la robotul din Unity la robotul real
		\item Al saptelea pas a constat în aplicarea valorilor primite asupra motoarelor și realizarea de modificări ulterioare.
	\end{itemize}
\end{comment}

\section(Arhitectura sistemului)








\printbibliography[heading=bibintoc,title=Bibliografie momentană]
\end{document}